\subsection{Evaluating a derivative on an integral lattice}
\hypertarget{source-integral-jets}{}

The coefficient estimates so far use products of framed minors. To pass to integral vectors of their rational span one needs saturation over the local valuation ring.

\begin{lemma}[Saturation of flag-minor polynomials]\label{lem:saturated-minor-polynomials}
Over a valuation ring \(\mathcal O\), fix a list of column
heights and take products of minors in the first \(h\) rows
at each height \(h\). Write \(S\) for their \(\mathcal O\)-span and \(A\) for the polynomial ring. For every \(f\in A\) and \(a\in\mathcal O\),
\[
 a\ne0,\quad af\in S\quad\Longrightarrow\quad f\in S.
\]

\end{lemma}

\begin{proof}
Sort the indices within each tableau column,
carrying the determinant sign; a repeated index kills the term.
One may straighten until the indices across each tableau row
are weakly increasing. Indeed suppose adjacent columns
\(a_1<\cdots<a_h\), \(b_1<\cdots<b_k\), \(k\le h\),
violate \(a_i\le b_i\). Antisymmetrize the two-minor product
in the \(h+1\) vectors indexed by
\(a_i,\ldots,a_h,b_1,\ldots,b_i\).

The antisymmetrized product vanishes because both minors use only the first \(h\) coordinates. Grouping shuffles
modulo permutations internal to the two lists gives an identity
with coefficients \(\pm1\) before collecting like terms. This
identity holds over the integers, as is seen by first expanding
over \(\mathbb Q\) and cancelling the common factorials.
With \(P\) the identity-shuffle product, the alternating relation has the form
\[
  P=-\sum_{\sigma\ne1}\operatorname{sgn}(\sigma)P_\sigma,
\]
where the sum uses the shuffle representatives and each \(P_\sigma\) retains the order inside the two lists. Every other shuffle
moves a value at most \(b_i<a_i\) into the left column in place
of a larger value. After sorting, a surviving term has a
strictly smaller left column in lexicographic order. Earlier
columns remain fixed. The whole tableau decreases in a finite
order, which proves spanning.

Order the matrix variables row by row, preferring smaller
column indices, and take lexicographic monomial order. The
leading monomial of a sorted minor is its diagonal product
with coefficient one. In a straightened tableau the row
multisets read off from the diagonal monomial recover each
weakly increasing row.
For a tableau with column entries \(a_{1c}<\cdots<a_{h_c c}\), the recovery rule is
\[
  \operatorname{in}(P)=\prod_c\prod_{i=1}^{h_c}X_{i,a_{ic}},\qquad \deg_{X_{ij}}\operatorname{in}(P)=\#\{c:a_{ic}=j\}.
\]
 Distinct such tableaux consequently have
distinct leading monomials. In an expansion over the fraction
field, the coefficient of the largest leading monomial is the
coefficient of its tableau. Subtract it and continue downwards;
if the polynomial is integral all coefficients so obtained
are integral. This proof works for products at several
letters because their matrices have disjoint variables.

\end{proof}

\begin{lemma}[Transfer of a coefficient bound to integral jet functionals]
\label{lem:jet-functional-coefficient-bound}
If \(\lambda\) is integral on the prescribed products of flag minors and \(P=\sum_\alpha c_\alpha X^\alpha\) is a local jet polynomial, then
\[
 v(\lambda(P))\ge\min_\alpha v(c_\alpha).
\]
 Joint integral remainders in independent
orbit variables obey the sum of the individual orbit bounds.
\end{lemma}

\begin{proof}
Write the jet polynomial in the frame \(Y_s\).
For a minor with row set \(I\) and column set \(J\),
Cauchy--Binet applies because the constant-coefficient
derivations commute:
\[
    \det((Y_s\partial)_{I,J})
      =\sum_{|K|=|I|}\det((Y_s)_{I,K})\det(\partial_{K,J}).
\]
The row set \(I\) is initial. After those derivatives act and
are evaluated at zero, only scalar coefficients of products
of initial-row minors remain. These are the polynomials of
Lemma~\ref{lem:saturated-minor-polynomials}. If the resulting
polynomial is nonzero with coefficient of least valuation \(a\),
divide by that coefficient to obtain an integral polynomial
in their fraction-field span. Saturation gives an integral
linear combination. Applying the integral functional to this
combination gives a value of valuation at least \(v(a)\)
after scaling back. The zero polynomial contributes zero.

For independent
sets of orbit variables an arbitrary integral factor \(U\) may
depend on all sets. Expand it in monomials before applying the
substitutions \(g\). Each monomial factors over the orbits, with
integral scalar, and only finitely many affect a specified jet
since \(g(0)=0\). Hence the bounds add even for this joint \(U\).
\end{proof}

There is a different calculation for a ramified quadratic factor.
At a place over \(p\), suppose the local series has the form
\[
 \exp\!\left(-S\sum_{t,a,b}u_t X_{tab}Y_{tab}\right) H(X,Y),
 \qquad v(S)\ge-1/2,\quad v(u_t)\ge0 ,
\]
and that every derivative of \(H\) at zero is integral. This last
property is preserved by integral linear substitutions, by the
chain rule with constant coefficients. Let \(H^\pm_t,A^\pm_t\)
be the total orders and numbers of minors on the two sides,
with \(X\) the minus side, and let \((A^-_t)'\) count minus
minors of height at least two.

For a fixed differential monomial
with respective entry exponents \(m^\pm_{tab}\), in the
derivative of the exponential times \(H\) each factor \(S\)
uses one differentiation from each side at a common entry.
At one entry write \(m^\pm=m^\pm_{tab}\). Choose \(h\) derivatives
from each side and match them by a bijection. The corresponding
coefficient and its valuation satisfy
\[
 C_h=\binom{m^-}{h}\binom{m^+}{h}h!(-Su_t)^h,
 \qquad v(C_h)\ge -h/2,\qquad 0\le h\le\min(m^-,m^+).
\]
The remaining derivatives act on \(H\) and have integral value. Summing
the losses over entries gives the lower bound
\(-\frac12\sum_{t,a,b}\min(m^-_{tab},m^+_{tab})\).

Evaluate the frames at independent indeterminates with their
Gauss valuation. The determinants of the square frames are
units there. Choose as the new minus derivations its given rows
and as the new plus derivations the inverse transpose rows.
For a change \(U\) of minus frame, the inverse transpose on the plus frame preserves the contraction because
\[
  (U x)^{\mathsf t}(U^{-\mathsf t}y)=x^{\mathsf t}y.
\]
 The plus frame is now an integral combination of the
new plus rows, and the minus minor operators use the identity
frame. Without restriction the sum of the minima at \(t\)
is at most \(H_t^-\).

At a selected \(t\in T\), the minus
first-row degree is exactly \(A^-_t\), while the plus first-row
degree is at most \(A^+_t\) after expanding its minors. Use
the plus count on that row and the minus counts on the other
rows. Writing \(a=1\) for the first row, this gives
\[
  \sum_{a,b}\min(m^-_{tab},m^+_{tab})\le \sum_b m^+_{t1b}+\sum_{a\ne1,b}m^-_{tab}\le A_t^+ +H_t^- - A_t^-.
\]
In particular the uniform lower bound used in the chain is
\[
 -\frac12\sum_t\left[H_t^-
       -{\bf1}_{t\in T}\{A^-_t-(A^-_t)'\}
       +{\bf1}_{t\in T} A^+_t\right].
\]
Passing back to the generic frame gives this bound coefficientwise.
An integral functional on its monomials preserves it. At a
place where all derivatives are integral and the differential
operator has integral coefficients, the bound is simply zero.

The calculations apply when the corresponding local form and
integral remainder have been obtained from the curve.
\label{cor:integral-jets-1}
